\begin{lemma}[Interchangeability across the two PE constituents]
Let $T$ be a positive-rung extended Tashkinov tree constructed from a
$k$-triple $G,e_0,\varphi$, with arbitrary certified input coloring,
terminal ladder tree $T_n$, terminal generating coloring $\varphi_n$,
and terminal stored-color set $D_n$.  Assume that $T$ has maximum
property and that its terminal extension is a PE.  Let $R_n$ and
$T_n\vee R_n$ be the closure and exact ordered union in a hierarchy
base under $\varphi_n$ (respectively `base.closure' and `base.core'
in Lean).

Assume, uniformly over every certified input coloring and generating
coloring, that every positive-rung ETT $T'$ of rung number $s<n$ with
maximum property satisfies all six lower-rung assertions:
\begin{enumerate}
\item $V(T')$ is elementary under its generating coloring;
\item under every proper coloring leaving exactly $e_0$ uncolored and
stable on its final ladder closure, that closure is interchangeable for
two colors whenever at least one of them is missing there;
\item stability descends at each PE support to the prefix before the
support and to its first-appearance ladder stage, and stored pairs at
distinct PE slots with distinct supports and the same first-appearance
index are disjoint;
\item at a terminal PE, the stored two-color component through the
support meets the terminal ladder tree only at that support;
\item every nonminimum endpoint of a defective boundary color precedes
each later PE support or SE extension vertex allowed by its
first-appearance index; and
\item every proper terminal stable coloring leaving exactly $e_0$
uncolored is a mod-coloring of the ETT.
\end{enumerate}

Let $\sigma_n$ be proper, leave exactly $e_0$ uncolored, and be
$T_n\oplus R_n,D_n,\varphi_n$-stable.

For arbitrary $\alpha,\beta\in[k]$, if
$\alpha\in\overline\sigma_n(T_n)$ and
$\beta\in\overline\sigma_n(R_n)$, then $\alpha,\beta$ are
$T_n\vee R_n$-interchangeable under $\sigma_n$.
All boundaries, trees, and two-color components use the original edge
type, so distinct parallel edges retain their identities.
\end{lemma}
\begin{proof}
Write $X=V(T_n)$, $Y=V(R_n)$, and
$U=V(T_n\vee R_n)$, with vertex sets understood as in
\noderef{TreeSequenceVertexSet}.  The generating-color field of
\noderef{ExtendedTashkinovTree} identifies $\varphi_n$ with the last
coloring of the certified series.  Eliminate the last certificate in
\noderef{TashkinovSeries} using the terminal-PE hypothesis.  The PE
constructor of \noderef{TashkinovSeriesStep} supplies the exact support
$v_n$, exact connecting-edge identity $f_n$, and colors
$\gamma_n,\delta_n$.  It also supplies properness and exact
one-uncoloredness of $\varphi_n$, and the equations recording the
terminal $\gamma_n,\delta_n$-switch.  No endpoint pair is substituted
for $f_n$.

The parallel clause of \noderef{HierarchyBase} says that $R_n$ is a
$\varphi_n$-closure of the ordered prefix $T_n(v_n)$, and supplies
an \noderef{OrderedTreeUnion} certificate with
\[
                         U=X\cup Y.                       \tag{1}
\]
It also says that the missing-color set of $T_n$ under the base
coloring is its missing-color set under $\varphi_n$.
The two conjuncts of \noderef{CombinedStableColoring}, with stability
expanded as in \noderef{StableColoring}, say that
$\sigma_n$ is $(T_n,D_n,\varphi_n)$-stable and simultaneously
$(R_n,\varnothing,\varphi_n)$-stable.  Applying
\noderef{CombinedStableClosurePreserved} gives, for every displayed edge
identity $g$ of either tree,
\[
             \sigma_n(g)=\varphi_n(g),                    \tag{2}
\]
and proves that the same ordered tree $R_n$, with the same edge
identities, is a $\sigma_n$-closure of $T_n(v_n)$.

We record three preliminary facts, with their hypotheses checked at the
coloring at which they will be used.

First, $U$ is elementary under $\sigma_n$.  Indeed, retain the same
core, closure, and support position from the hierarchy base.  Equality
of missing-color sets on $T_n$ is the vertex clause of the first stable
relation.  The preceding paragraph supplies the required
$\sigma_n$-closure, and the old ordered-union certificate is unchanged.
These data form a hierarchy base under $\sigma_n$.  The first stable
relation, properness, exact one-uncoloredness, positivity, maximum
property, the terminal-PE equation, and the displayed six lower-rung
assertions are precisely the hypotheses of
\noderef{ParallelClosureUnionElementary}.  Hence
\[
                  U\text{ is elementary under }\sigma_n.   \tag{3}
\]

By \noderef{InterchangePECommonMissingClosed}, a color missing
on both constituents is closed on $U$.  By
\noderef{InterchangePEClosureConstituent}, $R_n$ is interchangeable
whenever one of the two colors is missing there.  By
\noderef{InterchangePETreeConstituent}, the analogous statement holds
for $T_n$.  Finally, \noderef{InterchangePEClosedCore} gives
interchangeability on $U$ whenever a missing color there is closed.
These statements include singleton components through
\noderef{TwoColorPathComponent}, not only components containing an
edge-node.

It remains to prove (iv).  Choose
$u\in X$ missing $\alpha$ and $v\in Y$ missing $\beta$.
Suppose that roots $z_1,z_2$ witness two unequal full paper
$\alpha,\beta$-path components $C_1,C_2$ meeting $U$, possibly including
a singleton component, as in the expanded
definition \noderef{TreeInterchangeable}.  In particular,
\noderef{TwoColorPathComponent} for either root gives
$\alpha\ne\beta$.

We first locate the two missing-color witnesses.  If $u\in Y$, then
$\alpha$ is common missing, so (5) makes it closed on $U$; conclusion
(iii) would say that $\alpha,\beta$ are $U$-interchangeable, contrary to
$C_1,C_2$.  Hence $u\notin Y$.  If $v\in X$, then $\beta$ is common
missing and closed by (5).  Applying conclusion (iii) with the two colors
reversed gives $\beta,\alpha$ interchangeability.  The incidence graph,
path-component predicate, and \noderef{TreeInterchangeable} are symmetric
in their two displayed colors, so this again contradicts $C_1,C_2$.
Thus $v\notin X$, and
\[
                         u\in X-Y,\qquad v\in Y-X.             \tag{11}
\]
In particular $u\ne v$.  Elementariness (3), in the sense of
\noderef{ColoredVertexSetElementary}, together with $\beta$
missing at $v$ and $\alpha$ missing at $u$, shows respectively that
$\beta$ is present at $u$ and $\alpha$ is present at $v$.  Properness
then makes each of these present colors occur on one exact incident edge.
Consequently the full components $P_u,P_v$ through the vertex-nodes
$u,v$ each contain a genuine edge-node and have their root vertex of
degree one.  Along with $\alpha\ne\beta$, these give
\noderef{TwoColorNontrivialPathComponent}, and in particular
\noderef{TwoColorPathComponent}.

We now identify them with the counterexample components.  By (1), each
$C_i$ meets $X$ or $Y$.  Conclusion (i), applied with the common
$Y$-missing color $\beta$ and using color symmetry, permits at most one
of $C_1,C_2$ to meet $Y$.  Conclusion (ii), using $\alpha$ missing on
$X$, permits at most one to meet $X$.  Since both components meet
$X\cup Y$, no one of them can meet both constituents: the other component
would then have to meet one constituent and violate the corresponding
uniqueness.  After relabeling, $C_1$ meets $X$ and not $Y$, while $C_2$
meets $Y$ and not $X$.  Since $P_u$ is a path component meeting $X$,
conclusion (ii), applied to $P_u$ and $C_1$, identifies them.  Similarly
conclusion (i) identifies $P_v$ with $C_2$.  Thus $P_u\ne P_v$,
and \noderef{TwoColorComponentsDisjoint} makes their full node supports
disjoint.  If any vertex-node of $P_v$ lay in $X$, then $P_v$ and $P_u$
would be two path components meeting $T_n$, and conclusion (ii) would
force them equal.  Therefore the full component $P_v$ has no vertex in
$X$.

Define $\sigma'_n$ by \noderef{TwoColorComponentSwap}, swapping
$\alpha,\beta$ on the full component rooted at the vertex-node $v$.
To instantiate \noderef{OutsideKempeRecoloring}, take the singleton root
set $\{\operatorname{inl}v\}$.  The required distinct-color field is
$\alpha\ne\beta$ above.  The selected object is the full component
$P_v$, not merely a tail.  Its outside-$T_n$ field is exactly the absence
of all $X$-vertex-nodes just proved.  Finally, expanding
\noderef{TwoColorComponentSwap} says for each exact identity $f\in E$
that its value is transposed precisely when the edge-node
$\operatorname{inr}f$ belongs to $P_v$, and is fixed otherwise.  This is
word-for-word the edge-value field of
\noderef{OutsideKempeRecoloring} for the singleton root set.  Hence
$\sigma'_n$ is an outside Kempe recoloring of $T_n$.
Apply \noderef{OutsideKempeRecoloringStable} to $T_n$, with control
set $D_n$, the stated properness and one-uncoloredness, and this
outside-chain fact.  It proves that $\sigma'_n$ is proper, leaves
exactly $e_0$ uncolored, and is
$T_n,D_n,\sigma_n$-stable.  Composing with the given
$(T_n,D_n,\varphi_n)$-stability via \noderef{sc2} makes it
$(T_n,D_n,\varphi_n)$-stable.  This is exactly the stability conclusion
for which the manuscript invokes LEM-Stable; the direct
\noderef{OutsideKempeRecoloringStable} interface has exactly the hypotheses
available here.

We next preserve the part of the closure preceding $v$.  In the
augmentation sequence from $T_n(v_n)$ to $R_n$, consider an edge
adjoined before $v$.  Its color is missing at an earlier vertex.  By
(3) and \noderef{ColoredVertexSetElementary}, $\alpha$, missing at
$u\in X-Y$, is missing at no vertex of
$Y$; and $\beta$, missing at $v$, is missing at no earlier vertex
of $Y$.  Therefore none of these exact adjoined edge identities has
color $\alpha$ or $\beta$.  The swap on $P_v$ fixes every one of
them and preserves the missing-color witness at every preceding vertex.
Expanding \noderef{TashkinovAugmentation} at each step proves that the
same initial augmentation sequence is valid under $\sigma'_n$, and it
still reaches a tree containing $v$.

Let $S$ be the last tree of that copied finite sequence, so
$v\in V(S)$.  By \noderef{TashkinovClosureExists}, choose a
$\sigma'_n$-closure $R'_n$ of $S$.  Compose its augmentation sequence
with the copied sequence from $T_n(v_n)$ to $S$.  Its final closedness
is unchanged, so the definition of \noderef{TashkinovClosure} makes
$R'_n$ a closure of the original ordered prefix, and it still contains
$v$.

Construct the exact ordered union $T_n\vee R'_n$ by starting with all
entries of $T_n$ and then appending, in their $R'_n$-order, precisely
the entries whose new vertex is not already in $T_n$.  Every appended
edge attaches its new vertex to $T_n$ or to an earlier appended vertex,
because $R'_n$ is an augmentation sequence from a prefix of $T_n$.
Its new vertex and edge identity have not appeared earlier, by the
distinct-entry and attachment clauses of \noderef{TreeSequence} for
$T_n$ and $R'_n$.  Thus this is a
tree-sequence.  Directly from the construction, $T_n$ is its ordered
initial segment, its vertex set is the union, every new entry is the
same exact entry of $R'_n$, every required right-hand entry occurs, and
their relative order is retained.  These are all five fields of
\noderef{OrderedTreeUnion}.

The stable vertex clause identifies
the terminal missing-color set with that under $\varphi_n$; together
with the copied support, closure, and ordered union, this is a hierarchy
base under $\sigma'_n$.  Now all actual hypotheses of
\noderef{ParallelClosureUnionElementary} have been discharged: the
original ETT and terminal PE, positivity and maximum property, the full
arbitrary-input lower package, properness, exact one-uncoloredness,
$(T_n,D_n,\varphi_n)$-stability, and the newly constructed base.  Hence
$V(T_n\vee R'_n)$ is elementary under $\sigma'_n$.

The swap is disjoint from $T_n$, so $\alpha$ remains missing at
$u$.  At the end $v$ of $P_v$, swapping the component changes the
missing color from $\beta$ to $\alpha$.  The copied augmentation
sequence puts $v$ in $R'_n$.  Thus the two distinct vertices $u,v$
of $T_n\vee R'_n$ both miss $\alpha$ under $\sigma'_n$, contradicting
the elementarity just obtained by
\noderef{ColoredVertexSetElementary}.  Therefore two distinct components
cannot meet $U$, which is conclusion (iv).

Every use of a boundary, augmentation, component swap, or ordered union
above carries elements of the original edge type $E$; consequently no
step identifies distinct parallel edges having the same endpoints.
\end{proof}
